\begin{proposition}[E.29 --- Proved: a rational based $\mu_6$ action and three remaining coordinates]
\label{E-mu6-three-parameters}
Let $U=E\setminus E[3]$ be the cubic curve with the
specified tangential base point at $O$, and retain its
rational class-three Malcev group. The origin-fixing
geometric group $\mu_6$ has an actual based action on
this group, normalized by the arithmetic action:
\[
       g\beta g^{-1}=\beta^{\chi(g)\bmod6}
       \qquad(g\in G_{\mathbb Q}),
 \tag{E.29.1}
\]
where $\beta$ lifts a generator of $\mu_6$ and has
order six. Its cube is the rationalization of the
based involution of E.24. Thus the actual class
$[\eta]$ lies in the image of
\[
 H^1(G_{\mathbb Q,S_3},A_{\mathbb Q_3}^{\mu_6}(-1))
       \longrightarrow
 H^1(G_{\mathbb Q,S_3},A_{\mathbb Q_3}(-1)).
 \tag{E.29.2}
\]

Together with E.23, E.27 and E.28, this places the
class in the canonical rational direct sum
\[
 \operatorname{image}
 H^1(G_{\mathbb Q,S_3},B_{\mathbb Q_3}^{\mu_6}(-1))
 \ \oplus\
 H^1(G_{\mathbb Q,S_3},H(1)),
 \tag{E.29.3}
\]
of dimensions $2+1=3$. Here
$B_{\mathbb Q_3}^{\mu_6}(-1)\simeq\mathbb Q_3(1)^2$.
The first two coordinates are multiples of the
Kummer class of $3$; the third is the remaining
elliptic cohomology coordinate. Their actual values
are not determined by this proposition. The
denominator six in the based action and averaging
is used rationally, so no integral torsion coordinate
is discarded.
\end{proposition}
\begin{proof}
We construct the based action before quotienting by
inner path ambiguity. The origin-fixing automorphism
group of the $j=0$ elliptic curve is the group scheme
$\mu_6$ over $\mathbb Q$. It preserves $E[3]$.
Consequently the quotient stack
$\mathcal U=[U/\mu_6]$ is defined over $\mathbb Q$,
and $U\to\mathcal U$ is a finite etale torsor of
stacks. On geometric fundamental groups it gives
an extension by the six-element geometric group.

Near $O$, a formal parameter can be linearized for
this finite action in characteristic zero, keeping
its specified first-order tangent coordinate. To
see this explicitly, over a field containing
$\mu_6$ average any parameter $t_0$ with nonzero
derivative as
\[
       t=\frac16\sum_{j=0}^5
                 \zeta_6^{-j}\alpha^{j*}(t_0),
\]
after fixing the generator so that its derivative
is $\zeta_6$. Scaling the first coefficient to one
retains the given tangent. This is the eigenspace
projector, so the coordinate and its transformation
rule descend with the tangent character. Its
quotient coordinate is $q=t^6$. The map on nonzero
tangent directions sends the chosen direction
$t=1$ to $q=1$.

Let $d$ be the positively oriented geometric
meridian of this quotient punctured normal line,
and let $z_O$ be the meridian of the original cusp.
The actual covering $q=t^6$ gives
\[
                          d^6=z_O.
 \tag{E.29.4}
\]
The image of $d$ in the deck quotient is a generator
of $\mu_6$. Write $D$ for conjugation by $d$ on the
rational Malcev group of $N$. It preserves $z_O$
and satisfies $D^6=\operatorname{Int}(z_O)$.
The rational power $z_O^{-1/6}$ is defined by
$\exp(-\log(z_O)/6)$. Set
\[
                 \beta=\operatorname{Int}(z_O^{-1/6})D.
 \tag{E.29.5}
\]
Since $D$ fixes $z_O$, (E.29.4) gives $\beta^6=1$.
Its action on compact homology is the faithful
order-six elliptic automorphism, so its order is
exactly six.

Here is the arithmetic compatibility, including
the base-point comparison. The arithmetic action
on $N$ is the one from its specified tangent at
$O$. Sending that tangent through $q=t^6$ gives
the quotient tangent $1$, and the induced map of
root fibre functors identifies their arithmetic
splittings. Indeed $q^{1/n}$ pulls back to
$t^{6/n}$, and both fixed-root splittings fix these
compatible roots. A unit coordinate change with
constant coefficient one gives the same section.
Therefore, in this based orbifold extension, the
arithmetic section conjugates $d$ to $d^{\chi(g)}$
and $z_O$ to $z_O^{\chi_3(g)}$ in its rational
pro-three image. These equalities refer to the
actual tangential splitting, rather than an
arbitrary representative of the outer action.

Choose $r\in\{1,5\}$ congruent to $\chi(g)$
modulo six. In the local inertia extension,
\[
 \operatorname{Int}(d^{\chi(g)})
    =\operatorname{Int}
       \left(z_O^{(\chi_3(g)-r)/6}\right)D^r.
\]
This follows by subtracting the exponent $r$ in
the full local inertia group and using (E.29.4);
the remaining exponent is the stated three-adic
one after passage to the Malcev group. Combining
it with the conjugate of (E.29.5) gives
\[
 \begin{aligned}
 g\beta g^{-1}
 &=\operatorname{Int}(z_O^{-\chi_3(g)/6})
       \operatorname{Int}
          (z_O^{(\chi_3(g)-r)/6})D^r\\
 &=\operatorname{Int}(z_O^{-r/6})D^r
   =\beta^r.
 \end{aligned}
\]
This proves (E.29.1). In particular the action
commutes with $G_K$ for $K=\mathbb Q(\zeta_3)$,
and complex conjugation inverts its generator.
The cube of (E.29.5) is the order-two correction
for the same local normal-line extension. That
correction is unique because its rational inertia
has no two-torsion. It is therefore exactly the
rationalization of E.24's based involution.

The action on the associated graded module $A$
is the geometric $\mu_6$ action used in E.23:
inner conjugation acts trivially on degree three
modulo degree four. The line $\mathbb Q_3w$
is fixed by $\mu_6$, because the determinant of
the origin-fixing automorphism on compact
homology is one. Average any lift of $w$ in
$L_w\otimes\mathbb Q_3$ under the actual finite
group generated by $\beta$. Formula (E.29.1)
shows that this projector commutes with Galois.
The extension cocycle of the averaged lift
therefore has values in $A^{\mu_6}(-1)$.
Changing the lift changes only a coboundary,
proving (E.29.2).

We compute its kernel under projection to $Q=A/B$.
The geometric $\mu_6$ orbits on the nine cusps
are $\{O\}$, the two nonzero points killed by
$1-\zeta_3$, and the remaining six points.
The first subgroup is stable under conjugation
because $1-\zeta_3^2$ differs from $1-\zeta_3$
by a unit. All three orbits are consequently
Galois-stable. Their three orbit sums in cusp
homology have the single total-sum relation.
The compact summand has no $\mu_6$ invariants.
Thus
\[
 V_{\mathbb Q_3}^{\mu_6}
       =W^{\mu_6}\simeq\mathbb Q_3(1)^2,
 \qquad
 B_{\mathbb Q_3}^{\mu_6}(-1)
       \simeq V_{\mathbb Q_3}^{\mu_6}.
\]
This has an explicit basis in the notation of E.27.
On the Fermat cubic the origin is $O=C_0=[1:-1:0]$,
and the order-three automorphism
$[X:Y:Z]\mapsto[X:Y:\zeta_3 Z]$ fixes each $C_i$
and sends $A_i$ to $A_{i-1}$ and $B_i$ to $B_{i-1}$,
with subscripts modulo three. The involution interchanges
$A_i,B_i$ and interchanges $C_1,C_2$. Consequently
\[
 v_6=b_0+b_1+b_2,
       \qquad v_2=b_3
       \tag{E.29.6}
\]
are the orbit sums for the six-point and two-point orbits
and form a basis of $V_{\mathbb Q_3}^{\mu_6}$.
In the three Tate coordinates of E.27, ordered as
$b_0,b_1+b_2,b_3$, the two Kummer scalars attached
to the first two coordinates are therefore equal:
the vector of scalars has the form $(\lambda,\lambda,\mu)$.
This equality concerns the actual cohomology class under
the bracket map $v\mapsto[v,w]$ and the canonical
projector of E.27, because those maps commute with
the based $\mu_6$ action constructed above.
Rational invariants of a finite group are exact,
so $0\to B^{\mu_6}\to A^{\mu_6}\to Q^{\mu_6}
\to0$ is exact. The positive weights give
$H^0(Q^{\mu_6}(-1))=0$, and the $H^2$ vanishing
of E.16 gives $H^2(B^{\mu_6}(-1))=0$.
The corresponding $H^1$ sequence is therefore
short exact.

The full geometric averaging projector $e_{\mathcal G}$
on $A$, for $\mathcal G=E[3]\rtimes\mu_6$,
preserves $A^{\mu_6}$, kills $B$, and identifies
$A^{\mathcal G}$ with $Q^{\mathcal G}$, as in
E.27. It also preserves the compact/boundary
multigrading. E.28 kills the pure boundary
cubic component, so the allowed image in
$H^1(Q^{\mu_6}(-1))$ is precisely the embedded
$H^1(H(1))$ summand. The projector gives its
canonical lift in $H^1(A^{\mathcal G}(-1))$;
the complementary kernel is exactly the image
of $H^1(B^{\mu_6}(-1))$. This proves the direct
sum (E.29.3). Finally
$H^1(G_{\mathbb Q,S_3},\mathbb Q_3(1))$
is the one-dimensional rational $S_3$-unit
Kummer group, generated by $\kappa_3(3)$.
The two kernel coordinates and the one-dimensional
elliptic summand therefore have the stated
dimensions and meanings.
\end{proof}
